\section{Background and Related Work}
\label{sec:background}

Successive updates can partially undo one another while training loss continues to fall. We review the stability theory that explains this coexistence, the additional variability introduced by mini-batch sampling, methods that monitor instability during training, and work that connects update size to internal representations.

\subsection{Curvature, Oscillation, and the Edge of Stability}

Let $\theta_t\in\mathbb{R}^p$ collect the network's $p$ parameters after $t$ completed updates. The gradient $g_t$ is computed from a mini-batch at $t$, before the next update. Plain SGD uses a constant learning rate $\eta>0$ as
\begin{equation}
  \theta_{t+1}=\theta_t-\eta g_t.
  \label{eq:sgd_map_stability}
\end{equation}
For full-batch GD, $g_t=\nabla\mathcal L(\theta_t)$. The Hessian $H_t:=\nabla^2\mathcal L(\theta_t)$ describes how the slope of the loss changes with the parameters. On a fixed quadratic objective, its eigenvectors identify directions whose perturbations evolve independently: one update multiplies a perturbation along direction $i$ by $1-\eta\lambda_i$, where $\lambda_i$ is the curvature in that direction. If all curvatures are positive, every perturbation shrinks precisely when~\cite{polyak1964methods,bertsekas1997nonlinear}
\begin{equation}
  \max_i|1-\eta\lambda_i|<1
  \quad\Longleftrightarrow\quad
  \eta\lambda_{\max}(H_t)<2.
  \label{eq:quad_stability_condition}
\end{equation}
For $0<\eta\lambda_i<1$, a perturbation shrinks without changing sign; for $1<\eta\lambda_i<2$, it changes sign at every step but still shrinks; for $\eta\lambda_i>2$, it changes sign and grows. Back-and-forth motion can therefore be either convergent or unstable, even on a fixed quadratic.

In neural networks, $H_t$ changes as the parameters move, so Eq.~\eqref{eq:quad_stability_condition} gives a local stability threshold. Empirically, full-batch GD can drive the largest Hessian eigenvalue to or above $2/\eta$ even as training loss continues to fall over longer timescales~\cite{cohen2021edge,arora2022understanding}. Several accounts explain how training can sustain these oscillations. Self-stabilization describes how motion along a sharp direction can reduce curvature in that direction~\cite{damian2023self}. Central Flows describes the slower trajectory over which the oscillations occur~\cite{cohen2025central}. An analogous edge of stability has been observed for preconditioned sharpness under adaptive optimizers~\cite{cohen2022adaptive}.

These accounts leave open, however, how much of the distance traveled during oscillation remains as net parameter change. We measure this fraction by dividing the norm of the net displacement by the total length of the update path. In our GD experiments, this ratio drops sharply as consecutive updates begin to oppose one another. At the highest learning rates tested, only 2.6--5.4\% of the distance traveled remains as net displacement after 3,000 updates (Fig.~\ref{fig:gd-alignment-displacement}).

\subsection{Mini-Batch Sampling}

Different mini-batches produce different gradients at the same parameters~\cite{mandt2017sgd,hu2019diffusion,jastrzebski2018three}. Write $\xi_t$ for the difference between a batch gradient and the full gradient. For independent uniform batches of $b$ examples, each sampled without replacement from a dataset of $N>1$ examples,
\begin{equation}
\begin{aligned}
g_t&=\nabla\mathcal L(\theta_t)+\xi_t,\\
\mathbb E[\xi_t\mid\theta_t]&=0,\\
\operatorname{Cov}(\xi_t\mid\theta_t)
  &=\frac{N-b}{b(N-1)}\Sigma(\theta_t).
\end{aligned}
\label{eq:sgd_noise_decomposition}
\end{equation}
With per-example gradients $h_i:=\nabla\ell_i(\theta_t)$ and their mean $\mu:=N^{-1}\sum_i h_i=\nabla\mathcal L(\theta_t)$, the covariance
\begin{equation}
  \Sigma(\theta_t)
  :=\frac{1}{N}\sum_{i=1}^{N}(h_i-\mu)(h_i-\mu)^\top
  \label{eq:grad_cov_def}
\end{equation}
records how much individual examples disagree. Sampling therefore adds a source of gradient reversal that does not come from curvature. Edge of Stochastic Stability extends the stability analysis to mini-batch objectives and distinguishes sampling-induced oscillations from instability~\cite{andreyev2025stochastic}. We measure this effect through the cosine between consecutive batch gradients. In continuations from trained states, the smallest batches bring the negative mean cosine closer to zero in all three models. From initialization, even a positive full-batch mean can turn negative under sampling, as we observe in the MLP (Fig.~\ref{fig:sgd-moving-geometry}).

\subsection{Monitoring Curvature and Predicting Instability}

Neural-network Hessians are dominated by a few sharp directions~\cite{sagun2017hessian,ghorbani2019density}, so tracking $\lambda_{\max}$ is informative but requires repeated Hessian--vector products. Bai et al.~\cite{bai2026spikes} explain Adam spikes through decoupling between squared gradients and their running average, and identify this monitoring cost as a limitation. Barzilai--Borwein methods choose a step size from the parameter change and the gradient difference across a completed step~\cite{barzilai1988two,tan2016bbsgd}. We use these same quantities to recover the directional curvature averaged along the step the optimizer took. For full gradients of a smooth, fixed objective, this measurement is exact and equals the reciprocal of the BB1 step size when defined. It requires no additional backward pass or Hessian--vector product. Before the next Adam update, we use gradient reversal and norm growth together to warn of a loss increase of at least 0.1\%. Across all three models, the combined rule gives higher warning lift than either signal alone or a warning based on the preceding loss increase (Table~\ref{tab:severity}).

\subsection{From Loss to Layers and Update Control}

In practice, loss spikes are handled through global update size and training state: lowering the learning rate, tightening gradient clipping, adjusting Adam's second-moment decay, or rewinding to an earlier checkpoint and skipping batches~\cite{chowdhery2022palm,zhang2022opt,bai2026spikes}. These interventions change how training proceeds, but they leave open where in the network a loss excursion is accompanied by a reversal in representation. Central Flows leaves open how loss-landscape dynamics translate into the roles of individual layers~\cite{cohen2025central}. Representation-comparison methods characterize similarities across layers and training states~\cite{raghu2017svcca,kornblith2019similarity}; we instead track consecutive representation changes to locate reversal, quantify movement, and identify low-dimensional patterns and shared timing during oscillatory training. Kosson et al.~\cite{kosson2024warmup} identify representation change as a promising measure of update size, but leave its control open; their initial interventions focus on reducing warmup. We test local displacement caps during established oscillatory training. Because loss excursions can accompany improved feature learning~\cite{zhu2024catapults}, we assess training loss alongside fluctuation suppression.
